% Source: https://artofproblemsolving.com/community/c3h1429893_can_anyone_give_me_a_clue_hide_the_answer_please
% Stable identity: aops-topic-1429893-q01
We need $(10a+b)c=a(10b+c)$, that is $9ac+bc=10ab$, or
\[
9ac=b(10a-c).
\]
\textbf{$b$ not a multiple of $3$.} Then $9\mid10a-c$, so $c\equiv a\pmod9$, which for digits $1$--$9$ forces $c=a$; then $9a^2=9ab$ gives $b=a$, a trivial case.

\textbf{$b=3$.} $3ac=10a-c$, so $c(3a+1)=10a$. Since $\gcd(3a+1,a)=1$, $3a+1\mid10$, so $a=3$ and $c=3$: trivial.

\textbf{$b=6$.} $3ac=2(10a-c)$, so $c(3a+2)=20a$. Since $\gcd(3a+2,a)\mid2$, $3a+2$ divides $40$, so $3a+2\in\{5,8,20\}$, giving $a=1,c=4$ (the fraction $\dfrac{16}{64}$); $a=2,c=5$ ($\dfrac{26}{65}$); $a=6,c=6$ (trivial).

\textbf{$b=9$.} $ac=10a-c$, so $c(a+1)=10a$ and $a+1\mid10$: $a=1,c=5$ ($\dfrac{19}{95}$); $a=4,c=8$ ($\dfrac{49}{98}$); $a=9,c=9$ (trivial).

So the three other fractions are
\[
\boxed{\dfrac{19}{95}=\dfrac{1}{5},\qquad \dfrac{26}{65}=\dfrac{2}{5},\qquad \dfrac{49}{98}=\dfrac{4}{8}} .
\]
